\section{Four Levels of Formal Commitment}

Before introducing symbols, the chapter should state what sort of commitment
each symbol carries. Not every equation in a textbook has the same status.

\begin{enumerate}[nosep]
  \item \textbf{Descriptive notation.} Symbols that name analytically useful
        features, such as immediate reward pull or environmental reinforcement.
  \item \textbf{Heuristic scaffolds.} Simple summary expressions introduced for
        bookkeeping, ranking, or exposition.
  \item \textbf{Qualitative monotonic assumptions.} Claims such as ``more
        ambiguity usually raises transition difficulty'' that are plausible and
        structurally useful, but still not calibrated laws.
  \item \textbf{Unsupported if overstated.} Any claim of identified universal
        coefficients, uniquely correct functional form, or direct laboratory
        measurement for the chapter's everyday-action model.
\end{enumerate}

\begin{remark}
The practical rule for later chapters is simple: they may inherit the
notation, the scaffolds, and the monotonic assumptions, but they must not
quietly upgrade those elements into settled quantitative law.
\end{remark}
